\cxtitle{Exact QRCP can miss the orthonormal-row conditioning bound}

\begin{cxcredits}
\posedby{Anil Damle and Daniel Kressner \citeyearpar{AmselEtAl2026}}
\foundby{OpenAI Codex}{2026-08}
\formalizedby{OpenAI Codex}
\auditedby{OpenAI Codex}
\contributedby{Suvrit Sra}
\end{cxcredits}

\begin{cxcontext}
Let $Q\in\mathbb R^{n\times k}$ have orthonormal columns.  Exact
Businger--Golub QR with column pivoting (QRCP) applied to $Q^T$ repeatedly
chooses a remaining column with largest squared residual after orthogonal
projection away from the previously chosen columns.  If $P=QQ^T$ and the
previously chosen indices form $S$, the squared residual of index $i$ is the
Schur complement
\[
 \rho_S(i)=P_{ii}-P_{iS}P_{SS}^{-1}P_{Si}.
\]
Thus the exact pivot path depends only on the orthogonal projector $P$.
The bound quoted below asserts that some $I$ is well conditioned; what is at
issue is whether greedy pivoting finds one.  Randomized selection is the other
route: dual volume sampling applied to $Q^T$ draws $I$ with probability
proportional to $\det\bigl(Q(I,:)^TQ(I,:)\bigr)$ and admits polynomial-time
samplers \citep{LiJegelkaSra2017}.

Chen, Liu, He, and Dong first answered
Problem~\ref{ctx:qrcp-orthonormal-greedy} negatively by an asymptotic
low-coherence family \cite{ChenLiuHeDong2026}.  Their construction
prescribes a long upper-triangular selected block and then completes the
identity with several blocks of extra columns.  The finite witness below was
found later and independently.  It instead specifies a small rational point
in a Grassmann chart, takes its projector directly, and has only $k=3$ and
$n=8$.  No priority over Chen--Liu--He--Dong is claimed.
\end{cxcontext}

\begin{cxsource}{qrcp-orthonormal-greedy}
Problem~4.3 of Amsel et al. \cite{AmselEtAl2026}, sourced there to Anil Damle
and Daniel Kressner.  Chen, Liu, He, and Dong \cite{ChenLiuHeDong2026} were
the first to publish a negative resolution.
\end{cxsource}

\begin{cxstatement}{qrcp-orthonormal-greedy}
``It is known that there exists $I$ such that
$\|Q(I,:)^{-1}\|_2\leq\sqrt{k(n-k+1)}$ [HP92] and there is a fairly efficient
$O(nk^2)$ algorithm that satisfies this bound [Osi25].''  Problem~4.3 then
asks: ``Does the QRCP algorithm typically used in practice satisfy the above,
or a similar, bound.''
\end{cxstatement}

\begin{cxsummary}{qrcp-orthonormal-greedy}
For an exact rational rank-three projector on eight coordinates, QRCP has
strict pivots $I=(1,2,3)$ but the resulting inverse norm exceeds $\sqrt{18}$.
\end{cxsummary}

\begin{cxcertificate}{qrcp-orthonormal-greedy}
An exact rational $8\times3$ Grassmann-chart matrix, three positive rational
pivot gaps, and the Rayleigh gap $192051/250000$.
\end{cxcertificate}

\begin{cxrefutation}
Problem~\ref{ctx:qrcp-orthonormal-greedy} asks whether classical QRCP itself
attains the conditioning scale known to be achievable by a suitable
subset-selection algorithm \cite{AmselEtAl2026}.  The full question was first answered negatively by
Chen, Liu, He, and Dong \cite{ChenLiuHeDong2026}.  Here is a different,
small, exact counterexample to the displayed bound.

\begin{theorem}[\statusfalse: QRCP orthonormal-row bound]
\label{thm:qrcp-orthonormal-greedy}
There are $n=8$, $k=3$, and $Q\in\mathbb R^{8\times3}$ with $Q^TQ=I_3$
such that exact Businger--Golub QRCP applied to $Q^T$ makes three strict
pivot choices $I=(1,2,3)$, while
\[
 \|Q(I,:)^{-1}\|_2>\sqrt{k(n-k+1)}=\sqrt{18}.
\]
\end{theorem}
\begin{proof}
Put $V=\binom{I_3}{W}$, where
\[
 W=\frac1{1000}\begin{pmatrix}
 707&-23&-985\\
 1639&1259&-985\\
 1639&1259&-985\\
 -707&23&985\\
 1639&1259&-985
 \end{pmatrix},
 \qquad H=V^TV=I_3+W^TW,
\]
and define $Q=VH^{-1/2}$ using the positive-definite square root.  Then
$Q^TQ=I_3$, and its projector is the rational matrix
$P=QQ^T=VH^{-1}V^T$.  Hence all QRCP decisions can be certified using
rational arithmetic alone.

For completeness, exact Schur-complement evaluation gives the following
pivot residual and runner-up residual at each step:
\[
\begin{array}{c|c|c|c|c}
 S & \text{pivot} & \rho_S(\text{pivot}) &
 \max_{i\notin S\cup\{\text{pivot}\}}\rho_S(i)&\text{gap}\\ \hline
 \varnothing&1&\dfrac{100874772187}{197275840682}&
 \dfrac{99899751883}{197275840682}&
 \dfrac{487510152}{98637920341}\\[5pt]
 \{1\}&2&\dfrac{4179375000}{14410681741}&
 \dfrac{28781505000}{100874772187}&
 \dfrac{474120000}{100874772187}\\[5pt]
 \{1,2\}&3&\dfrac{8000}{46809}&
 \dfrac{38809}{234045}&
 \dfrac{397}{78015}
\end{array}
\]
Every gap is positive, so tie handling is irrelevant and QRCP selects
$I=(1,2,3)$.

The selected matrix is $Q(I,:)=H^{-1/2}$, so
\[
 \|Q(I,:)^{-1}\|_2^2=\lambda_{\max}(H)
 =1+\lambda_{\max}(W^TW).
\]
For the integer vector $x=(2,1,-1)^T$, direct exact arithmetic gives
\[
 \|Wx\|_2^2-17\|x\|_2^2=\frac{192051}{250000}>0.
\]
Consequently $\lambda_{\max}(W^TW)>17$, whence
$\|Q(I,:)^{-1}\|_2^2>18=k(n-k+1)$.  This is the claimed strict
counterexample.
\end{proof}

\begin{remark}[Scope]
Problem~\ref{ctx:qrcp-orthonormal-greedy} asks whether QRCP satisfies ``the
above, or a similar, bound''.
What falls here is the first half: the displayed bound $\sqrt{k(n-k+1)}$ is
violated outright at $n=8$, $k=3$.  A single finite witness says nothing about
the second half, since a ``similar'' bound loose enough to admit this instance
is not excluded by it.  The asymptotic family of Chen, Liu, He, and Dong
\cite{ChenLiuHeDong2026} is what settles the broader reading; this witness is
smaller and exactly rational, not stronger.
\end{remark}
\end{cxrefutation}
